\begin{textsample}{POS Dim 1 – 2019\_c5 – Score 18.00 – 2019\_c5/enade\_2019\_eng\_prod\_c7c0e9871a.txt}  \label{ex:f1_pos_047}
CAPÍTULO 5 DISTRIBUIÇÃO DOS CONCEITOS 5.1 \textbf{PANORAMA} NACIONAL DA DISTRIBUIÇÃO DOS CONCEITOS O Conceito Enade é um indicador de qualidade que avalia os cursos por \textbf{intermédio} dos desempenhos dos estudantes a partir dos resultados do Enade.
Os valores possíveis vão de 1 ( pior situação ) a 5 ( melhor situação ).
Os cursos com apenas um ou sem nenhum concluinte participante não obtêm o Conceito Enade, ficando Sem Conceito ( SC ).
Na Tabela 5.1, são apresentadas a quantidade e a distribuição de cursos de Engenharia de Produção participantes do Enade/2019, por faixa de conceito e Grande Região.
Enfatiza -se, mais uma vez, que a diferença entre os cursos \textbf{tabulados} neste capítulo e no capítulo 2 corresponde aos cursos que não foram avaliados, ou seja, àqueles, em princípio, sem alunos concluintes que tivessem sido inscritos para a prova.
Observando -se os dados da Tabela 5.1, nota -se que, dos 654 cursos participantes, 249 ( 38,1 \% ) classificaram -se com conceito 2, o valor modal.
Esse foi também o conceito modal em quase todas as regiões : Norte ( 58,3 \% ), Nordeste ( 38,6 \% ), Sudeste ( 39,6 \% ) e Centro-Oeste ( 40,6 \% ).
O conceito 3 foi o segundo mais frequente em nível nacional ( 36,9 \%, correspondendo a 241 cursos ), sendo também o conceito modal da região Sul ( 46,2 \% ).
O conceito 4 foi o terceiro mais frequente com 81 cursos, 12,4 \% dos cursos foram avaliados com esse conceito.
Houve, ainda, 40 cursos ( 6,1 \% ) que receberam conceito 5, e 23 cursos ( 3,5 \% ) que receberam conceito 1.
Dos 654 cursos de Engenharia de Produção, 20 ( 3,1 \% ) ficaram Sem Conceito ( SC ).
A região Norte participou com 24 cursos ou 3,7 \% do total nacional.
Desses, 14 cursos receberam o conceito 2, o que \textbf{equivale} a 58,3 \% do total regional, além de ser o conceito modal na região, como já comentado.
A cinco cursos ( 20,8 \% ) atribuiu -se conceito 3, a dois cursos ( 8,3 \% ), conceito 4 e a um curso ( 4,2 \% ), conceito 5.
Dois cursos ficaram Sem Conceito ( SC ) e nenhum curso recebeu o conceito 1 nessa região.
A região Nordeste participou com 101 cursos ou 15,4 \% do total nacional.
Nessa região, 39 cursos ( 38,6 \% em termos regionais ) obtiveram conceito 2, o conceito modal para a região, como já comentado.
Os conceitos 1, 3 e 4 foram atribuídos, respectivamente, a sete, 35 e a 13 cursos ( 6,9 \%, 34,7 \% e 12,9 \%, respectivamente ).
Ainda nessa região, cinco cursos ( 5,0 \% ) foram avaliados com o conceito 5.
Dois cursos ficaram Sem Conceito ( SC ).
Dos 367 cursos participantes da região Sudeste ( 56,1 \% do total nacional ), o conceito 2 foi o conceito modal, como já anteriormente assinalado, tendo sido obtido por 144 cursos ( 39,2 \% ).
O conceito 1 foi atribuído a dez cursos ( 2,7 \% ), e o conceito 3, a 132 cursos ( 36,0 \% ).
Receberam os conceitos 4 e 5, respectivamente, 45 cursos ( 12,3 \% ) e 25 cursos ( 6,8 \% ).
Nessa região, onze cursos ficaram Sem Conceito ( SC ).
Os 130 cursos da região Sul corresponderam a 19,9 \% do total nacional.
Como já \textbf{indicado}, houve predominância do conceito 3, atribuído a 46,2 \% dos cursos da região ( conceito modal ), ou seja, atribuído a 60 dos 130 cursos participantes na região Sul.
O conceito 2 foi atribuído a 39 cursos ( 30,0 \% ), o conceito 4, a 18 cursos ( 13,8 \% ), o conceito 5, a sete cursos ( 5,4 \% ) e o conceito 1, a dois cursos ( 1,5 \% ).
Quatro cursos ficaram Sem Conceito ( SC ) na região Sul.
Já dos 32 cursos participantes na região Centro-Oeste ( 4,9 \% do total nacional ), 13 ( 40,6 \% em termos regionais ) receberam conceito 2, o conceito modal, como já fora destacado.
Nove cursos ( 28,1 \% ) obtiveram o conceito 3, quatro cursos ( 12,5 \% ), o conceito 1, três cursos ( 9,4 \% ), o conceito 4 e dois cursos ( 6,3 \% ), o conceito 5.
Além disso, nessa região, um curso ficou Sem Conceito ( SC ).
No Gráfico 5.1, é apresentada a distribuição \textbf{cumulativa} do Conceito Enade segundo Grande Região.
Em um gráfico de distribuição \textbf{cumulativa}, em cada valor de \textbf{abscissa} ( eixo x ), o valor apresentado no eixo das \textbf{ordenadas} ( eixo y ) é o valor acumulado da \textbf{variável} ( conceito Enade ).
Observando a linha vertical \textbf{pontilhada}, que ajuda a \textbf{visualizar} o percentual de conceitos acumulados até o valor 3, ou seja, conceitos 1, 2 e 3, pode -se notar que, por exemplo, a região Norte ( linha azul ) apresenta mais de 85 \% dos cursos com conceito 3 ou abaixo deste valor.
As demais regiões apresentam entre 80 \% e 84 \% dos cursos na mesma situação : com conceito 3 ou abaixo.
Destaca -se que em todas as regiões o conceito 5 foi alcançado por algum curso, e as cinco \textbf{poligonais} só alcançam 100 \% no conceito máximo.
Em linhas gerais, em um gráfico de distribuição \textbf{cumulativa}, \textbf{poligonais} mais à esquerda representam Grandes Regiões com pior distribuição, e \textbf{poligonais} mais à direita, Grandes Regiões com melhores conceitos.
A região Norte ( linha azul ) apresenta os cursos com conceitos mais baixos, ao passo que as regiões Sudeste ( linha vermelha ) e Sul ( linha \textbf{rosa} ) apresentam as distribuições com valores maiores.
Cursos em IES situadas nas regiões Nordeste ( linha verde ) e Centro-Oeste ( linha laranja ) apresentam uma situação intermediária. 5.2 CONCEITOS POR CATEGORIA ADMINISTRATIVA, POR MODALIDADE DE ENSINO E GRANDE REGIÃO Na Tabela 5.2, é apresentada a distribuição dos cursos de Engenharia de Produção participantes do Enade/2019, por Categoria Administrativa e Modalidade de ensino, de acordo com os conceitos por eles alcançados, segundo as Grandes Regiões brasileiras.
Dos 654 cursos participantes, 114 ( 17,4 \% ) eram \textbf{ministrados} em Instituições \textbf{Públicas}, e 540 ( 82,6 \% ), em Instituições Privadas.
Quanto à Modalidade de ensino, 23 cursos eram a Distância, e 631, Presenciais.
De acordo com as informações da Tabela 5.2, em termos nacionais, entre Instituições \textbf{Públicas} e Privadas, observa -se que, dos 40 cursos avaliados com conceito 5, 33 eram oferecidos em IES Públicas, e sete, em IES Privadas.
Dos 114 cursos participantes de IES Públicas, o conceito 4 foi o valor modal, atribuído a 41 cursos ( 36,0 \% ).
Dos demais cursos de Instituições \textbf{Públicas} participantes, um curso ( 0,9 \% ) foi avaliado com conceito 1, sete cursos ( 6,1 \% ) foram avaliados com conceito 2, 32 cursos ( 28,1 \% ) foram avaliados com conceito 3 e, como já comentado, 33 cursos ( 28,9 \% ), com conceito 5.
Nesta categoria, nenhum curso ficou Sem Conceito ( SC ).
Na rede privada, o conceito modal foi 2, atribuído a 242 cursos, o correspondente a 44,8 \% dos 540 cursos da categoria.
Dos demais cursos participantes, 22 ( 4,1 \% ) receberam conceito 1, e 209 ( 38,7 \% ), conceito 3.
O conceito 4 foi atribuído a 40 cursos ( 7,4 \% ), e o conceito 5, a sete cursos ( 1,3 \% ).
Nesta Categoria Administrativa, 20 ( 3,7 \% ) cursos ficaram Sem Conceito ( SC ).
Dos 631 cursos participantes da modalidade presencial, a distribuição foi \textbf{bimodal}, com os conceitos 2 e 3 concentrando 236 cursos ( 37,4 \% ) cada.
Dos demais cursos nesta modalidade de ensino, 22 cursos ( 3,5 \% ) receberam conceito 1, 81 cursos ( 12,8 \% ), conceito 4, e 38 cursos ( 6,0 \% ), conceito 5.
Nesta modalidade de ensino, 18 cursos ( 2,9 \% ) ficaram Sem Conceito ( SC ).
Aos 23 cursos a Distância informados atribuíram -se os conceitos 1 ( um curso, 4,3 \% ), 2 ( 13 cursos, 56,5 \% ), 3 ( cinco cursos, 21,7 \% ) e 5 ( dois cursos, 9,5 \% ).
Nessa Modalidade de ensino, dois cursos ficaram Sem Conceito ( SC ) e nenhum recebeu o conceito 4.
Na análise por região, observa -se que, na região Norte, as Instituições Privadas participaram com 18 cursos ( 75,0 \% do total regional ), dos quais dois ficaram Sem Conceito ( SC ), e nenhum recebeu os conceitos 1 ou 5.
O conceito modal para as Instituições Privadas na região Norte foi 2, atribuído a 13 cursos, correspondendo a 72,2 \% dessa categoria na região.
Dois cursos foram avaliados com conceito 3, e um curso, com conceito 4.
As Instituições Públicas participaram com seis cursos na região Norte ( 25,0 \% do total regional ), aos quais foram atribuídos os conceitos 2, 4 e 5 ( com um curso cada ) e o conceito 3 ( três cursos, modal ).
Nenhum curso ficou Sem Conceito ( SC ) ou com conceito 1 nessa combinação de categoria e região.
Todos os cursos oferecidos na região Norte eram presenciais, com conceito modal 2 ( 14 cursos ).
Na região Nordeste, a rede privada concentrou 79 dos 101 cursos participantes, o equivalente a 78,2 \% do total da região.
Dentre os cursos oferecidos por IES Privadas no Nordeste, 38 cursos foram avaliados com conceito 2, o conceito modal.
Sete cursos obtiveram conceito 1, 28 cursos receberam conceito 3 e quatro cursos receberam conceito 4.
Nessa combinação de categoria e região, dois cursos ficaram Sem Conceito ( SC ) e nenhum recebeu conceito 5.
As Instituições \textbf{Públicas} da região Nordeste participaram com 22 cursos ( 21,8 \% do total da região ).
Desses, nove obtiveram conceito 4, o valor modal, seguidos por outros sete cursos com conceito 3, cinco cursos com conceito 5 e um curso com conceito 2.
Nenhum curso ficou Sem Conceito ( SC ) ou com conceito 1 nessa região.
A quase totalidade dos cursos oferecidos na região Nordeste eram presenciais, dos 100 cursos presenciais 39 cursos receberam o conceito 2 ( modal ), dois cursos ficaram Sem Conceito ( SC ).
Os demais cursos foram alocados nos conceitos 1 ( sete cursos ), 3 ( 34 cursos ), 4 ( 13 cursos ) e 5 ( cinco cursos ).
Um curso foi oferecido na modalidade a Distância na região Nordeste que foi avaliado com o conceito 3.
Na região Sudeste, a proporção de cursos da rede privada, 85,3 \%, foi mais elevada do que a registrada nas demais regiões brasileiras, correspondendo a 313 dos 367 cursos participantes.
Nessa categoria e região, o conceito modal foi 2, atribuído a 142 cursos e, onze cursos ficaram Sem Conceito ( SC ).
Os demais foram avaliados com conceito 1 ( dez cursos ), conceito 3 ( 121 cursos ), conceito 4 ( 23 cursos ) e conceito 5 ( seis cursos ).
Dos 54 cursos oferecidos por Instituições \textbf{Públicas} na região Sudeste, o conceito modal foi 4, atribuído a 22 cursos.
Nessa região, dois cursos de IES Públicas receberam conceito 2, 11 cursos o conceito 3 e 19 cursos o conceito 5.
Também na região Sudeste, a maioria dos cursos foi oferecido na modalidade Presencial : 350 cursos.
Para esses, a moda foi o conceito 2, com 134 cursos.
Os demais cursos presenciais da região foram alocados nos conceitos 1 ( dez cursos ), 3 ( 129 cursos ), 4 ( 45 cursos ) e 5 ( 23 cursos ).
Ficaram Sem Conceito ( SC ) nove cursos.
Na região Sudeste, os 17 cursos oferecidos na modalidade Educação a Distância obtiveram os conceitos 2 ( dez cursos, modal ), 3 ( três cursos ) e 5 ( dois cursos ).
Dois cursos ficaram Sem Conceito ( SC ).
As Instituições Privadas concentraram 106 dos 130 cursos participantes da região Sul, 81,5 \% do total regional.
Desses, 51 cursos obtiveram conceito 3, o conceito modal.
Nesta combinação de Categoria Administrativa e Grande Região, dois cursos receberam conceito 1, 37 cursos, conceito 2, 11 cursos, conceito 4 e, ainda, um curso, conceito 5.
Quatro cursos ficaram Sem Conceito ( SC ).
As Instituições \textbf{Públicas} na região Sul participaram com 24 cursos ( 18,5 \% dos cursos da região ), aos quais foram atribuídos os conceitos 2 ( dois cursos ), 3 ( nove cursos, modal ), 4 ( sete cursos ) e 5 ( seis cursos ).
Nenhum curso da região Sul oferecido por IES Pública ficou Sem Conceito ( SC ) ou recebeu o conceito 1.
A maioria dos cursos oferecidos na região Sul eram presenciais, ou seja, 126 cursos dos 130 presentes na região.
O conceito modal foi o 3 com 59 cursos, e o conceito 5 concentrou sete cursos.
Os demais cursos foram concentrados nos conceitos 2 ( 36 cursos ), 4 ( 18 cursos ) e 1 ( dois cursos ).
Ainda, outros quatro cursos ficaram Sem Conceito ( SC ) na modalidade Presencial na região Sul.
Dos quatro cursos da modalidade a Distância, três cursos receberam o conceito 2 e um curso o conceito 3.
Na região Centro-Oeste, 24 dos 32 cursos participantes eram de Instituições Privadas ( 75,0 \% em termos regionais ).
Desses, 12 receberam conceito 2, o conceito modal.
Dos demais cursos, sete receberam conceito 3, três, conceito 1 e um recebeu conceito 4.
Um curso ficou Sem Conceito ( SC ), e nenhum curso recebeu conceito 5.
Quanto aos oito cursos oferecidos por Instituições \textbf{Públicas} na região Centro-Oeste ( 25,0 \% do total regional ) todos foram conceituados.
Os conceitos 3, 4 e 5 receberam dois cursos cada e os conceitos 1 e 2 receberam um curso cada.
A quase totalidade dos cursos oferecidos na região Centro-Oeste eram presenciais.
Dos 31 cursos da modalidade Presencial da região Centro-Oeste, um curso ficou Sem Conceito ( SC ).
Os demais foram distribuídos aos conceitos 1 ( três cursos ), 2 ( 13 cursos ), 3 ( nove cursos ), 4 ( três cursos ) e 5 ( dois cursos ).
O curso da modalidade a Distância dessa região recebeu o conceito 1.
No Gráfico 5.2, apresenta -se a distribuição do Conceito Enade, segundo a Categoria Administrativa da IES.
Os cursos em IES Públicas ( linha azul ) apresentam uma \textbf{poligonal} mais à direita do que os em IES Privadas ( linha verde ) e, consequentemente, uma distribuição de Conceitos Enade com valores maiores. 1, 37 cursos, conceito 2, 11 cursos, conceito 4 e, ainda, um curso, conceito 5.
Quatro cursos ficaram Sem Conceito ( SC ).
As Instituições \textbf{Públicas} na região Sul participaram com 24 cursos ( 18,5 \% dos cursos da região ), aos quais foram atribuídos os conceitos 2 ( dois cursos ), 3 ( nove cursos, modal ), 4 ( sete cursos ) e 5 ( seis cursos ).
Nenhum curso da região Sul oferecido por IES Pública ficou Sem Conceito ( SC ) ou recebeu o conceito 1.
A maioria dos cursos oferecidos na região Sul eram presenciais, ou seja, 126 cursos dos 130 presentes na região.
O conceito modal foi o 3 com 59 cursos, e o conceito 5 concentrou sete cursos.
Os demais cursos foram concentrados nos conceitos 2 ( 36 cursos ), 4 ( 18 cursos ) e 1 ( dois cursos ).
Ainda, outros quatro cursos ficaram Sem Conceito ( SC ) na modalidade Presencial na região Sul.
Dos quatro cursos da modalidade a Distância, três cursos receberam o conceito 2 e um curso o conceito 3.
Na região Centro-Oeste, 24 dos 32 cursos participantes eram de Instituições Privadas ( 75,0 \% em termos regionais ).
Desses, 12 receberam conceito 2, o conceito modal.
Dos demais cursos, sete receberam conceito 3, três, conceito 1 e um recebeu conceito 4.
Um curso ficou Sem Conceito ( SC ), e nenhum curso recebeu conceito 5.
Quanto aos oito cursos oferecidos por Instituições \textbf{Públicas} na região Centro-Oeste ( 25,0 \% do total regional ) todos foram conceituados.
Os conceitos 3, 4 e 5 receberam dois cursos cada e os conceitos 1 e 2 receberam um curso cada.
A quase totalidade dos cursos oferecidos na região Centro-Oeste eram presenciais.
Dos 31 cursos da modalidade Presencial da região Centro-Oeste, um curso ficou Sem Conceito ( SC ).
Os demais foram distribuídos aos conceitos 1 ( três cursos ), 2 ( 13 cursos ), 3 ( nove cursos ), 4 ( três cursos ) e 5 ( dois cursos ).
O curso da modalidade a Distância dessa região recebeu o conceito 1.
No Gráfico 5.2, apresenta -se a distribuição do Conceito Enade, segundo a Categoria Administrativa da IES.
Os cursos em IES Públicas ( linha azul ) apresentam uma \textbf{poligonal} mais à direita do que os em IES Privadas ( linha verde ) e, consequentemente, uma distribuição de Conceitos Enade com valores maiores. 5.3 CONCEITOS POR ORGANIZAÇÃO ACADÊMICA E GRANDE REGIÃO Na Tabela 5.3, encontra -se a distribuição dos conceitos atribuídos aos cursos participantes do Enade/2019, na Área de Engenharia de Produção, por Organização Acadêmica, segundo as Grandes Regiões brasileiras.
Dos 654 cursos de Engenharia de Produção participantes, 242 eram oferecidos em Universidades, 178 em Centros Universitários, 223 em Faculdades e 11 em CEFET/IF.
Esta distribuição corresponde a, respectivamente, 37,0 \%, 27,2 \%, 34,1 \% e 1,7 \% dos cursos.
De acordo com os dados apresentados, de todos os 40 cursos avaliados com conceito 5, 32 eram vinculados a Universidades.
Os cursos vinculados a esse tipo de Organização Acadêmica tiveram o conceito 3 como conceito modal, atribuído a 99 cursos ( 40,9 \% ).
Os demais cursos vinculados a Universidades receberam os conceitos 1 ( quatro cursos ), 2 ( 53 cursos ), 4 ( 52 cursos ) e conceito 5 ( 32 cursos, como já mencionado ).
Ficaram Sem Conceito ( SC ), dois cursos.
Dos 178 cursos em Centros Universitários, o conceito modal foi 2, atribuído a 81 cursos ( 45,5 \% ).
Quatro cursos vinculados a esse tipo de Organização Acadêmica ficaram Sem Conceito ( SC ), e os outros cursos receberam os conceitos 1 ( três cursos ), 3 ( 78 cursos ), 4 ( dez cursos ) e 5 ( dois cursos ).
Dos 223 cursos mantidos por Faculdades, 115 ( 51,6 \% ) receberam conceito 2, o conceito modal.
O segundo conceito mais frequentemente atribuído aos cursos oferecidos por essa categoria de IES foi o conceito 3, recebido por 61 cursos.
Na sequência, 16 cursos obtiveram conceito 1, 14 cursos, conceito 4, e três cursos, conceito 5.
Dos cursos oferecidos em Faculdades, 14 ficaram Sem Conceito ( SC ).
Dos 11 cursos mantidos por CEFET/IF, cinco ( 45,5 \% ) receberam conceito 4, o conceito modal.
Os conceitos 3 e 5 foram atribuídos a três cursos cada.
Nenhum curso ficou Sem Conceito ( SC ) ou recebeu conceito 1 e 2.
Considerando -se, separadamente, as regiões brasileiras, verifica -se que, na o Norte, as Universidades participaram com apenas sete dos 24 cursos avaliados.
O conceito modal das Universidades da região Norte também foi o 3 com três cursos, o conceito 2 foi recebido por dois cursos e os conceitos 4 e 5 por um curso cada.
Nessa combinação de organização e região, nenhum curso ficou Sem Conceito ( SC ) ou foi avaliado com o conceito 1.
Os Centros Universitários da região Norte foram representados também por sete cursos, destes quatro cursos obtiveram conceito 2.
Um curso ficou Sem Conceito ( SC ), outro recebeu o conceito 3 e um outro o conceito 4.
Nenhum curso recebeu os conceitos 1 ou 5, nessa combinação de categoria e região.
As Faculdades participaram com dez cursos na região Norte, os quais receberam os conceitos 2 ( oito cursos, modal ) e 3 ( um curso ).
Nenhum curso oferecido em Faculdades da região obteve conceito 1, 4 ou 5.
Um curso ficou Sem Conceito ( SC ).
Os CEFET/IF não participaram com cursos na região Norte.
Na região Nordeste, as Universidades participaram com 31 dos 101 cursos da Área de Engenharia de Produção oferecidos.
Nenhum desses cursos ficou Sem Conceito ( SC ) ou recebeu conceito 1.
O conceito modal 3 foi atribuído a 12 cursos.
Os cursos \textbf{restantes} obtiveram conceitos 2 ( cinco cursos ), 4 ( dez cursos ) e 5 ( quatro cursos ).
Os Centros Universitários \textbf{contaram} com 28 cursos participantes na região Nordeste, 15 dos quais receberam o conceito modal 3.
Para os demais cursos, foram atribuídos os conceitos 2 ( 11 cursos ) e 4 ( dois cursos ).
Nenhum curso oferecido por Centros Universitários no Nordeste ficou Sem Conceito ( SC ) ou recebeu os conceitos 1 ou 5.
As Faculdades foram representadas por 41 cursos na região Nordeste, 23 dos quais receberam conceito 2, o conceito modal.
Oito cursos receberam o conceito 3, sete cursos receberam o conceito 1 e o conceito 4 recebeu um curso.
Na região Nordeste, dois cursos mantidos por Faculdades ficaram Sem Conceito ( SC ) e nenhum o conceito 5.
Os CEFET/IF participaram com um curso na região Nordeste avaliado com conceito 5.
Na região Sudeste, as Universidades concentraram 136 dos 367 cursos de Engenharia de Produção da região.
Dos cursos oferecidos em Universidades, nessa região, o conceito modal foi 3, atribuído a 52 cursos.
Os demais cursos receberam os conceitos 1 ( dois cursos ), 2 ( 36 cursos ), 4 ( 26 cursos ) e 5 ( 18 cursos ).
Dois cursos ficaram Sem Conceito ( SC ).
Os Centros Universitários participaram com 104 cursos na região Sudeste, dos quais 49 obtiveram o conceito modal 2.
O conceito 1 foi atribuído a três cursos, e o conceito 3, a 42 cursos.
Três cursos ficaram Sem Conceito ( SC ) e dois receberam o conceito 5 nessa combinação de organização acadêmica e região.
Já as Faculdades foram representadas por 117 cursos na região Sudeste, dos quais seis ficaram Sem Conceito ( SC ).
O conceito modal 2 foi atribuído a 59 cursos.
Aos demais cursos foram atribuídos os conceitos 1 ( cinco cursos ), 3 ( 35 cursos ), 4 ( nove cursos ) e 5 ( três cursos ).
Os CEFET/IF participaram com dez cursos que foram avaliados com os conceitos 3 ( três cursos ), 4 ( cinco cursos, modal ) e 5 ( dois cursos ).
Dos 130 cursos da região Sul, 56 cursos eram mantidos por Universidades, para os quais o conceito modal foi 3, atribuído a 29 cursos.
Os demais cursos receberam os conceitos 2 ( oito cursos ), 4 ( 12 cursos ) e 5 ( sete cursos ).
Nesse tipo de organização, nenhum curso da região Sul ficou Sem Conceito ( SC ) ou recebeu conceito 1.
Dos 35 cursos participantes de Centros Universitários da região Sul, nenhum ficou Sem Conceito ( SC ) ou recebeu os conceitos 1 e 5.
Os conceitos \textbf{restantes}, 2, 3 e 4, foram atribuídos a 15, 18 ( modal ) e dois cursos, respectivamente.
Quanto aos 39 cursos vinculados a Faculdades na região Sul, 16 receberam o conceito modal 2.
O conceito 1 foi atribuído a dois cursos, o conceito 3, a 13 cursos, e o conceito 4, a quatro cursos.
Nenhum curso recebeu conceito 5 e quatro cursos ficaram Sem Conceito ( SC ).
Os CEFET/IF não participaram com cursos na região Sul.
Na região Centro-Oeste, 12 dos 32 cursos oferecidos eram mantidos por Universidades.
Para este tipo de Organização Acadêmica, três cursos foram avaliados com os conceitos 3 e 4 cada ( \textbf{bimodal} ) e dois cursos foram avaliados com os conceitos 1, 2 e 5 cada.
Nesse tipo de organização, nenhum curso ficou Sem Conceito ( SC ).
Os Centros Universitários da região Centro-Oeste \textbf{contaram} com quatro cursos, avaliados com conceitos 2 e 3 ( dois cursos cada ).
Dos 16 cursos oferecidos em Faculdades na região Centro-Oeste, nove receberam o conceito modal 2, dois obtiveram conceito 1, quatro obtiveram conceito 3.
Um curso ficou Sem Conceito ( SC ) e nenhum curso recebeu os conceitos 4 e 5.
Os CEFET/IF não participaram com cursos na região Centro-Oeste.
No Gráfico 5.3, apresenta -se a distribuição do Conceito Enade dos cursos de Engenharia de Produção, segundo a Organização Acadêmica da IES.
Os cursos em CEFET/IF ( linha laranja ) apresentam uma \textbf{poligonal} mais à direita do que as demais, \textbf{denotando} uma melhor distribuição de Conceitos Enade.
No outro extremo, com distribuição muito semelhante entre si, encontram -se os cursos mantidos por Centros Universitários ( linha verde ) e Faculdades ( linha vermelha ), apresentando as respectivas \textbf{poligonais} mais à esquerda e, consequentemente, distribuição de Conceitos Enade com menores valores.
Em posição intermediária encontram -se os cursos mantidos por Universidades ( linha azul ).

% matched lemmas: abscissa, bimodal, contar, cumulativo, denotar, equivaler, indicado, intermédio, ministrar, ordenado, panorama, poligonal, pontilhado, pública, restante, rosa, tabular, variável, visualizar
\end{textsample}
